%\input{../preamble}
%\usepackage{nomencl}
%
%\newcommand{\ra}{\mathfrak{r}}
%\DeclareMathOperator{\Fix}{{\bf{Fix}}}
%\DeclareMathOperator{\MEAS}{{\bf{MEAS}}}
%\newcommand{\G}{\mathcal{G}}
%\DeclareMathOperator{\MAlg}{{\bf{MAlg}}}
%\newcommand{\tr}{\operatorname{tr}}
%\begin{document}
%
%One common way to construct equivalence relations and groupoids is by considering either orbit equivalence relations or transformation groupoids induced by group actions.
In this section we will study actions which induce sofic transformation groupoids. Again, we will focus only on $\ra$-discrete, standard probability-measure preserving groupoids.

%Moreover, we will only consider group actions which allow us to construct probability measure-preserving groupoids, which we define below. (These are the same actions as considered in Proposition \ref{propositionactionsmeasurepreserving}.)

\begin{definition}\label{definitionpmpaction}
Let $G$ be a group and $\theta$ and $(X,\mu)$ a measure space. An action $\theta$ of $G$ by Borel automorphisms of $X$ is \emph{measure-preserving} if $\mu(\theta_g(A))=\mu(A)$ for all Borel $A\subseteq X$ and $g\in G$.% (or equivalently if the transformation groupoid $G\ltimes_\theta X$ is probability measure-preserving, when endowed with the measure $\mu$).
\end{definition}

\begin{proposition}\label{propositionactionsmeasurepreserving}
Let $\theta$ be an action of a countable group $G$ on a standard probability space $(X,\mu)$ by Borel automorphisms of $X$. Then the following are equivalent:
\begin{enumerate}
\item[(1)] $\theta$ measure-preserving.% is invariant by each automorphism $\theta_g$, $g\in G$ (that is, the push forward measure $(\theta_g)_*\mu$ coincides with $\mu$ for all $g\in G$);
\item[(2)] $\mu$ is invariant for the Borel groupoid $G\ltimes_\theta X$;
\item[(3)] $\mu$ is invariant for the orbit equivalence relation $\mathcal{R}(\theta)$.% (Example \ref{exampleorbitgroupoid})
\end{enumerate}
\end{proposition}
\begin{proof}
The equivalence (2)$\iff$(3) is proven in Proposition \ref{propositionmeasureinvariancefororbitrelation}.

For every $g\in G$ and Borel $A\subseteq X$, the element $\left\{g\right\}\times A$ of $\BOR(G\ltimes_\theta A)$ satisfies
\[\so(\left\{g\right\}\times A)=A,\qquad\ra(\left\{g\right\}\times A)=\theta_g(A)\]
and every element of $\BOR(G\rtimes_\theta A)$ is in fact a countable disjoint union of these elements. The equivalence (1)$\iff$(2) follows immediately.\qedhere
\end{proof}

Recall that recall that equivalence relations correspond to principal groupoids (Proposition \ref{principalequalequivalencerelation}), and that \emph{free} actions are those which induce principal transformation groupoids (Example \ref{examplewhentransformationgroupoidisprincipal}). Since we are working in the category of measured spaces, null sets can be essentially disregarded.%However we are considering measure spaces and measure-preserving actions, so we will essentially disregard null sets (i.e., of measure zero), as usual.

\begin{definition}
A probability measure-preserving action $\theta$ of a countable group $G$ on a measure space $(X,\mu)$ is (\emph{essentially}) \emph{free} if for all $g\in G$, the set $\Fix(\theta_g)=\left\{x\in X:\theta_g(x)=x\right\}$ of points fixed by $g$ is null.
\end{definition}

%. Again, when dealing with probability spaces we are lenient with null sets.

\begin{definition}
A standard, $\ra$-discrete probability measure-preserving groupoid $(\G,\mu)$ is (\emph{essentially}) \emph{principal} if a.e.\ point of $\G[0]$ has trivial isotropy, i.e., if
\[\mu\left(\left\{x\in\G[0]:\G_x^x=\left\{x\right\}\right\}\right)=1.\]
\end{definition}

\begin{remark}
If $\theta$ is an essentially free probability measure-preserving action of a countable group $G$ on a standard probability space $(X,\mu)$, we can find a conull subset $Y$ of $X$ such that $\theta_g(Y)\subseteq Y$ for all $g\in G$, and $\theta_g(x)\neq x$ for all $x\in Y$, namely,
\[Y=\bigcap_{g\in G}\supp(\theta_g)=\left\{x\in X:\forall g\in G,\ \theta_g(x)\neq x\right\}.\]
A similar comment holds for for $\ra$-discrete, standard probability measure-preserving groupoids. Therefore our usage of ``free'' and ``principal'', in the category of measure spaces, will not be contradictory with the usual one, in the category of sets.

Moreover, certain classes of topological groupoids are called ``essentially principal'', (\cite{MR584266,MR2876963,MR3377390}). To avoid any confusion, we will use ``topologically principal'' and ``effective'' instead for topological groupoids (see Definitions \ref{definitiontopologicallyprincipal} and \ref{definitioneffective}).
\end{remark}

The following can be proven with the same arguments of Example \ref{examplewhentransformationgroupoidisprincipal} and the remark above.

\begin{proposition}
A probability measure-preserving action $\theta$ of a countable group $G$ on a probability space $(X,\mu)$ is free if and only if the transformation groupoid $G\ltimes X$ is principal.
\end{proposition}
%\begin{proof}
%As usual, we identify $(G\ltimes X)^{(0)}$ with $X$. For all $x\in X$, we can identify $(G\ltimes X)_x^x$ with the isotropy subgroup $G_x=\left\{g\in G:\theta_g(x)\right\}$, via the map
%\[G_x\to (G\ltimes X)_x^x,\qquad g\mapsto (g,x)\]
%and thus
%\[\left\{x\in X: (G\ltimes X)_x^x\neq\left\{x\right\}\right\}=\bigcup_{g\neq 1}\Fix(\theta_g)\supp\]
%and since $G$ is countable, the left-hand side is null if and only if $\Fix(\theta_g)$ is null for all $g\neq 1$, which is the equivalence we wanted.\qedhere
%\end{proof}

We will now define \emph{Bernoulli shifts}. Let $G$ be any countable group and $X$ any set. We consider the action $L$ of left multiplication of $G$ on itself: $L_g:G\to G$, $L_g(h)=gh$ for all $h,g\in G$, and then dualize $L$ to an action $\beta$ on the space $X^G$ of functions from $G$ to $X$ by setting
\[\beta_g(f)=f\circ L_{g^{-1}}\qquad\text{ for all}\quad f\in X^G\quad\text{and}\quad g\in G.\]
If we use the product notation $X^G=\prod_G X=\left\{(x_h)_{h\in G}:x_h\in X\right\}$ instead, the action $\beta$ becomes
\[\beta_g\left((x_h)_{h\in G}\right)=\left(x_{g^{-1}h}\right)_{h\in G}\]
for all $g\in G$ and all $(x_h)_{h\in G}\in X^G$.

\begin{definition}
The action $\beta$ above is called the \emph{Bernoulli shift} of $G$ on $X$.
\end{definition}

Bernoulli shifts are one of the basic examples in Symbolic Dynamics (see \cite{MR1222644}). In fact, they can be used to model other actions, as the simple proposition below shows. (Recall Definition \ref{definitioncovariant} for a \emph{covariant} map.)

\begin{proposition}
Let $\theta$ be any action of a countable group $G$ on a set $X$, and let $\beta$ be the Bernoulli shift of $G$ and $X$. Then the map $F:X\to X^G$, $F(x)=(\theta_{g^{-1}}(x))_{g\in G}$ is injective and covariant.
\end{proposition}

\begin{proposition}
Let $G$ be a nontrivial countable group and $(X,\mu)$ a standard probability space. Endow $X^G$ with the product measure coming from $\mu$. Then the Bernoulli shift $\beta$ of $G$ on $X$ is a probability measure-preserving action of $G$, and it is free if and only if $(X,\mu)$ is non-atomic, or if $G$ is infinite and $(X,\mu)$ is not an atom.
\end{proposition}
\begin{proof}
The invariance of measure is immediate on cylinder sets (those of the form $\prod_{g\in G} X_g$, where $X_g\subseteq X$ is Borel), and since these generate the Borel structure of $X^G$ then $\beta$ is measure-preserving.
\begin{itemize}
\item First assume that $(X,\mu)$ is non-atomic. Then $(X,\mu)$ it is isomorphic to the interval $[0,1]$ with Lebesgue measure (\cite[Theorem 17.41]{MR1321597}), and
\[\Fix(\theta_g)=\left\{(x_h)_h\in X^G:(x_{g^{-1}h})_h=(x_h)_h\right\}\subseteq\left\{(x_h)_h\in X^G:x_1=x_g\right\}\]
The latter set can be regarded as $\Delta\times X^{G\setminus\left\{1,g^{-1}\right\}}$, where $\Delta$ is the diagonal of $X\times X$, and so it has measure $0$.

\item Now assume that $G$ is infinite and $X$ is not an atom. Choose distinct elements $h_1,h_2,\ldots,\in G$ such that all $h_1,h_2,\ldots,g^{-1}h_1,g^{-1}h_2,\ldots$ are pairwise distinct. Then
\[\Fix(\beta_g)\subseteq\left\{(x_h)_h\in X^G:x_{g^{-1}h_i}=x_{h_i}\text{ for all }i\right\},\]
and similarly to the previous case above we can see the latter set as an infinite product of diagonals in $X\times X$ and copies of $X$. The diagonal in $X\times X$ has measure smaller than $1$ (because $X$ is not an atom), so an infinite product of them will have measure $0$.
\end{itemize}

If $(X,\mu)$ has an atom $\left\{x_0\right\}$ and $G$ is finite, then $\left\{(x_0)_h\right\}$ is an atom of $X^G$ and $\beta_g((x_0)_h))=(x_0)_h$ for all $g\in G$. If $X=\left\{x_0\right\}$ then $X^G=\left\{(x_0)_h\right\}$ is also an atom. In any of these cases, $\beta$ is not free.\qedhere
\end{proof}

\begin{remark}
In fact, the classification, up to covariant isomorphisms, of Bernoulli shifts has led to the the development of \emph{Entropy Theory} for measure-preserving actions of groups. Initially, entropy theory was developed for actions of $\mathbb{Z}$, then extended to actions of amenable groups, and recently a theory of entropy for sofic groups has been developed by Bowen \cite{MR2552252}, and subsequently improved by Kerr and Li \cite{MR2854085}. We refer to the survey paper \cite{MR3666024} for more details and the history of entropy theory.
\end{remark}

We will now be interested in sofic actions.

\begin{definition}\index{Sofic!Action}
A probability measure-preserving action $\theta$ of a countable group $G$ on a space $(X,\mu)$ is \emph{sofic} if the probability measure-preserving groupoid $(G\ltimes X,\mu)$ is sofic.
\end{definition}

For the remainder of this section, let us recall and fix some notation.

\subsubsection*{Conventions}\label{conventionssoficactions}

\begin{itemize}
\item If $X$ is a set and $f\in\mathcal{I}(X)$, we regard $\Fix(f)$ and $\supp(f)$ as in Examples \ref{examplefixedidempotentsinix} and \ref{examplesupportix}, i.e.,
\[\Fix(f)=\left\{x\in\dom(f):f(x)=x\right\},\qquad\supp(f)=\left\{x\in\dom(f):f(x)\neq x\right\}.\]
\item We will consider only countable groups, standard probability spaces and measure-preserving actions.
\item We again consider finite sets $Y_n$ with $n$ elements, endowed with their normalized counting measures $\mu_\#$, the permutation group $\mathfrak{S}_n$ of $Y_n$, the measure algebra $\MAlg(Y_n)$ and the measured semigroups $\MEAS(Y_n^2)=\mathcal{I}(Y_n)$ endowed with their normalized Hamming distances $d_\#$ and traces $\tr_\#$. We will also use the symbols $\mu_\#$, $d_\#$ and $\tr_\#$ for corresponding functions on ultraproducts.
\item $\mathcal{U}$ will always denote a ultrafilter of sets on $\mathbb{N}$.
\end{itemize}
Let $\theta$ be a measure-preserving action of $G$ on a standard probability space $(X,\mu)$.
\begin{itemize}
\item We identify each $g\in G$ with the element $\left\{g\right\}\times X$ of the full group $[G\ltimes_\theta X]$ (compare with Example \ref{exampleginsidebisectionsoftransformationgroupoid}). We also identify $E(\MEAS(G\ltimes X))$ with  $\MAlg(X)$, so if $g\in G$ and $A\in\MAlg(X)$, the product $gA$ corresponds to the element $\left\{g\right\}\times A$ of $\MEAS(G\ltimes X)$.
\item We consider the actions of $G$ on $\MAlg(X)$ and of $\prod_\mathcal{U}\mathfrak{S}_{n_k}$ on $\prod_\mathcal{U}\MAlg(Y_{n_k})$ as in Definition \ref{definitionactionofinvertibleones}:
\[g\cdot A=\theta_g(A),\qquad(\sigma_k)_\mathcal{U}\cdot(A_k)_\mathcal{U}=(\sigma_k(A_K))_\mathcal{U}.\]
\end{itemize}

We can compare actions of countable groups with the notion of \emph{weak containment} (an analogous notion for representations has been widely studied; see \cite{MR2415834}): By Proposition \ref{propositionultraproductoffunctions}, every probability measure-preserving action $\theta$ of $G$ on a probability space $(Y,\nu)$ induces an action $\theta^\mathcal{U}$ of $G$ on $\prod_\mathcal{U}\MAlg(Y)$ by
\[\theta^\mathcal{U}_g(A_n)_\mathcal{U}=(\theta_g(A_n))_\mathcal{U}.\]

%and it is routine to check that $\theta^\mathcal{U}$ is an action of $G$.)\nomenclature[Z]{$\theta^\mathcal{U}$}{Induced action on ultraproduct.}

\begin{definition}[{\cite[p.\ 64]{MR2583950}}]\index{Weak containment}
Let $\alpha$ and $\theta$ be free probability measure-preserving actions of a countable group $G$ on spaces $(X,\mu)$ and $(Y,\nu)$, respectively. We say that $\alpha$ is \emph{weakly contained} in $\theta$ if there exists an isometric morphism
\[\phi:\MAlg(X)\to\prod_\mathcal{U}\MAlg(Y)\]
which is covariant with respect to $\alpha$ and $\theta^\mathcal{U}$.
\end{definition}

In other words, if an action $\alpha$ is weakly contained in $\theta$ then $\alpha$ can be modelled, in the sense of continuous logic, by $\theta$. Another important property of Bernoulli shifts, initially proved by Abért and Weiss in \cite{MR3035287} and later extended by Seward in \cite{arxiv1501.03367v3}, is that Bernoulli shifts are weakly contained in every other free probability measure-preserving action.

\begin{theorem}[{\cite[Theorem 1.1]{MR3035287}}]\label{theoremabertweiss}
Let $\beta$ be the Bernoulli shift of a countable group $G$ on a standard probability space $(X,\mu)$ (where $X^G$ is endowed with the product measure). Then any free probability-measure preserving action of $G$ on any standard probability space $(Y,\nu)$ weakly contains $\beta$.
\end{theorem}

We will now describe soficity of an action in terms of covariant morphisms.

\begin{lemma}\label{lemmaseparatesetwithdisjointimages}
Suppose $X$ and $Y$ are Borel spaces with $Y$ standard, and $f,g:X\to Y$ are Borel maps with $f(x)\neq g(x)$ for all $x\in X$. Then there exists a countable partition $\left\{A_n\right\}$ of $X$ such that $f(A_n)\cap g(A_n)=\varnothing$ for all $n$.

Consequently, if $\mu$ is a probability measure on $X$, then for all $\epsilon>0$ there exists a finite Borel partition $A_1,\ldots,A_n,C$ of $X$ with $f(A_i)\cap g(A_i)=\varnothing$ for $1\leq i\leq n$ and $\mu(C)<\epsilon$.
\end{lemma}
\begin{proof}
Let $d$ be a metric on $Y$ compatible with the Borel structure and for which $(Y,d)$ is separable. Let $I_1=\left\{x\in X:d(f(x),g(x))>1\right\}$ and for $k\geq 2$, consider the Borel set
\[I_k=\left\{x\in X:1/k<d(f(x),g(x))\leq 1/(k-1)\right\}\]
so $\left\{I_k\right\}_k$ is a Borel partition of $X$. Given $k$, let $\left\{B^k_n\right\}_n$ be a Borel partition of $Y$ by sets of diameter at most $1/k$, and let $A^k_n=f^{-1}(B^k_n)\cap I_k$. If $x\in A^k_n$, then $f(x)\in B^k_n$ and $d(g(x),f(x))>1/k\geq \operatorname{diam}B^k_n$, so $g(x)\not\in B^k_n$. This proves
\[f(A^k_n)\cap g(A^k_n)\subseteq B^k_n\cap g(A^k_n)=\varnothing\qquad\text{for all}\quad k\text{ and }n,\]
and $\left\{A^k_n:k,n\right\}$ is a countable Borel partition of $X$.\qedhere
\end{proof}

\begin{lemma}\label{lemmatraceissmallerundercovariant}
Let $\theta$ be a probability measure-preserving action of $G$ on $(X,\mu)$. Suppose $\sigma:G\to\prod_\mathcal{U}\mathfrak{S}_n$ and $\phi:\MAlg(X)\to\prod_\mathcal{U}\MAlg(Y_{n_k})$ are covariant semigroup morphisms, with $\phi$ isometric. Then for all $g\in G$ and $A\in\MAlg(X)$, $\tr(\sigma(g)\phi(A))\leq\tr(gA)$.
\end{lemma}
\begin{proof}
Given $g\in G$ and $A\in\MAlg(X)$, let $B=\operatorname{supp}\theta_g\cap A$. Given $\epsilon>0$, choose, by Lemma \ref{lemmaseparatesetwithdisjointimages}, a finite partition $B_1,\ldots,B_n,C$ of $B$ with $\theta_g(B_i)\cap B_i=\varnothing$ for $1\leq i\leq n$ and $\mu(C)<\epsilon$.

Then the elements $\phi(B_i)$ are pairwise disjoint, contained in $\phi(A)$, and covariance of $(\sigma,\phi)$ implies
\[\left[\sigma(g)\cdot\phi(B_i)\right]\land\phi(B_i)=\phi(g\cdot B_i)\cap\phi(B_i)=\phi(\theta_g(B_i)\cap B_i)=0,\]
and from this we obtain $\phi(B_i)\subseteq\supp\left(\sigma(g)\right)\land \phi(A)=\supp(\sigma(g)\phi(A))$. It follows that
\begin{align*}
\operatorname{tr}(\sigma(g)\phi(A))&=\mu_\#(\phi(A))-\mu_\#(\supp(\theta_g\phi(A)))\leq\mu_\#(\phi(A))-\sum_{i=1}^n\mu_\#(\phi(B_i))\\
&=\mu(A)-\sum_{i=1}^n\mu(B_i)=\tr(gA)+\mu(C)\leq\tr(gA)+\epsilon.
\end{align*}
Since $\epsilon$ is arbitrary, we obtain the result.\qedhere
\end{proof}

The next theorem formalizes the description of soficity used by Ozawa in \cite{ozawasoficnotes}.

\begin{theorem}\label{theoremrewriesoficactionintermsofcovariant}
A free probability measure-preserving action $\theta$ of $G$ on $(X,\mu)$ is sofic if and only if there are a group morphism $\sigma:G\to\prod_\mathcal{U}\mathfrak{S}_n$ and an isometric morphism $\phi:\MAlg(X)\to\prod_\mathcal{U}\MAlg(Y_{n_k})$ such that $(\sigma,\phi)$ is covariant with respect to the canonical actions.
\end{theorem}

Recall, from the conventions in page \pageref{conventionssoficactions}, that we identify $G$ and $\MAlg(X)$ with subsemigroups of the measured semigroup $\MEAS(G\ltimes_\theta X)$. The proof of this theorem is similar in parts to that of Theorem \ref{theoremfirstimportant}.

\begin{proof}[Proof of \ref{theoremrewriesoficactionintermsofcovariant}]
First, suppose $\theta$ is sofic, and let $\Phi:\MEAS(G\ltimes_\theta X)\to\prod_\mathcal{U}\MEAS(Y_{n_k}^2)$ be a sofic embedding. Restricting $\Phi$ to $G$ and to $\MAlg(X)$ we obtain the morphisms $\sigma$ and $\phi$, and $\phi$ is isometric. If $A\in\MAlg(X)$ and $g\in G$,
\[\phi(g\cdot A)=\Phi(gAg^{-1})=\sigma(g)\phi(A)\Phi\sigma(g)^{-1}=\sigma(g)\cdot\phi(A),\]
so $(\sigma,\phi)$ is covariant.

Conversely, suppose the covariant pair $(\sigma,\phi)$ exists. We need a few steps in order to define a sofic embedding of $G\ltimes_\theta X$.

\begin{enumerate}
\item {\bf{If $g,h\in G$ and $A,B\in\MAlg(X,\mu)$ satisfy $A\cap B=\theta_g(A)\cap\theta_h(B)=\varnothing$, then $\sigma(g)\phi(A)$ and $\sigma(h)\phi(B)$ are compatible in $\prod_\mathcal{U}\MEAS(Y_{n_k}^2)$.}}% (see Definition \ref{definitioncompatibility}).

By covariance, we have
\[\left(\sigma(g)\phi(A)\right)^*\left(\sigma(h)\phi(B)\right)=\phi(A)\sigma(g^{-1}h)\phi(B)=\phi(A\cap(g^{-1}h\cdot B))\sigma(g^{-1}h)\]
and since
\[A\cap(g^{-1}h\cdot B)=\theta_{g^{-1}}(\theta_g(A)\cap\theta_h(B))=\varnothing,\]
then $\left(\sigma(g)\phi(A)\right)^*\left(\sigma(h)\phi(B)\right)=0$, and similarly $\left(\sigma(g)\phi(A)\right)*\left(\sigma(h)\phi(B)\right)^*=0$, so $\sigma(g)\phi(A)$ and $\sigma(h)\phi(B)$ are compatible.

\item {\bf{If $A_1,\ldots,A_n$ are pairwise disjoint and $g\in G$, then $\sigma(g)\phi(\bigcup_{i=1}^n A_i)=\bigvee_{i=1}^n\sigma(g)\phi(A_i)$}}.

First note that $\bigvee_{i=1}^n\sigma(g)\phi(A_i)$ is well-defined by item 1., so this follows from distributivity of $\prod_\mathcal{U}\MEAS(Y_{n_k}^2)$ and the fact that $\phi$ preserves joins (Theorem \ref{theoremtracepreservingisisometric}).
\end{enumerate}

Every element $B$ of $\MEAS(G\ltimes_\theta X)$ can be written uniquely as a countable join $B=\bigcup_{g\in G}gB_g$, where $B_g:=\so(B\cap(\left\{g\right\}\times X))$. Let $\mathfrak{M}$ be the submonoid of $\MEAS(G\ltimes_\theta X)$ consisting of those $B\in\MEAS(G\ltimes_\theta X)$ for which all except finitely many $B_g$ are empty. Define
\[\Theta:\mathfrak{M}\to\prod_\mathcal{U}\MEAS(Y_{n_k}^2),\qquad\Theta(B)=\bigvee_{g\in G}\sigma(g)\phi(B_g),\]
which is well-defined by item 1.\ above. To verify that $\Theta$ is a morphism, note that
\[AB=\bigcup_{g,h\in G}gA_ghB_h=\bigcup_{g,h\in G}gh((h^{-1}\cdot A_g)\cap B_h)=\bigcup_{k\in G}k\left(\bigcup_{g\in G}((g^{-1}k)^{-1}\cdot A_g)\cap B_{g^{-1}k}\right),\]
and similarly $\Theta(A)\Theta(B)=\bigvee_{k\in G}\sigma(k)\left(\bigvee_{g\in G}((\sigma(g)^{-1}\sigma(k))^{-1}\cdot\phi(A_g))\land\phi(B_{g^{-1}k})\right)$, so the facts that $(\sigma,\phi)$ is covariant and item 2.\ above imply that $\Theta(AB)=\Theta(A)\Theta(B)$.
%\[\Theta(AB)=\bigvee_{k\in G}\sigma(k)\left(\bigvee_{g\in G}(\sigma(g^{-1}k)^{-1}\cdot\phi(A))\land \phi(B_{g^{-1}k})\right)=\bigvee_{g,h\in G}\sigma(g)\sigma(h)\]
 Let us prove that $\Theta$ is trace-preserving.

Since $\theta$ is free then $\tr(gB_g)=0$ whenever $g\neq 1$, and Lemma \ref{lemmatraceissmallerundercovariant} implies also that $\tr(\sigma(g)\phi(B_g))=0$ when $g\neq 1$. Therefore
\[\tr(\Phi(B))=\sum_{g\in G}\tr(\sigma(g)\phi(B_g))=\mu(\phi(B_1))=\mu(B_1)=\sum_{g\in G}\tr(gB_g)=\tr(B).\]

Therefore $\Theta$ is trace-preserving, and hence isometric on the dense submonoid $\mathfrak{M}$ of $\MEAS(G\ltimes_\theta X)$, and thus extends uniquely to a sofic embedding of $G\ltimes_\theta X$.\qedhere
\end{proof}

\begin{definition}\index{Covariant sofic embedding}\index{Sofic embedding!covariant}
A covariant pair $(\sigma,\phi)$ as in the Theorem \ref{theoremrewriesoficactionintermsofcovariant} will be called a \emph{covariant sofic embedding}.
\end{definition}

In terms of approximation properties, as in \ref{theoremequivalenceembeddingintoultraproductformetricstructures}, this theorem can be rewritten as follows:

\begin{corollary}\label{corollaryalmostcovariant}
A free probability measure-preserving action $\theta$ of $G$ on $(X,\mu)$ is sofic if and only if for every $\epsilon>0$ and for every pair of finite subsets $F\subseteq G$ and $\mathcal{A}\subseteq\MAlg(X)$, there are $n\in\mathbb{N}$ and maps $\sigma:G\to\mathfrak{S}_n$ and $\phi:\MAlg(X)\to\MAlg(Y_n)$ satisfying
\begin{enumerate}
\item[(i)] $d_{\#}(\sigma(gh),\sigma(g)\sigma(h))<\epsilon$;
\item[(ii)] $d_{\#}(\phi(A\cap B),\phi(A)\cap\phi(B))<\epsilon$;
\item[(iii)] $|\mu_\#(\phi(A))-\mu(A)|<\epsilon$;
\item[(iv)] $d_{\#}(\phi(\theta_g(A)),\sigma(g)\cdot\phi(A))<\epsilon$
\end{enumerate}
for all $g,h\in F$ and all $A,B\in\mathcal{A}$.
\end{corollary}

\begin{definition}\index{Almost covariant morphism}
A pair of maps $(\sigma,\phi)$ satisfying properties (i)-(iv) above is called an \emph{$(F,\mathcal{A},\epsilon)$-almost covariant morphism}.
\end{definition}

As an immediate consequence of this description of soficity, we obtain:

\begin{corollary}
Soficity of free actions is preserved under weak containment: If $\alpha$ is a free sofic action which weakly contains a free action $\beta$, then $\beta$ is sofic as well.
\end{corollary}

Therefore, if an infinite countable group $G$ admits a free sofic action then any Bernoulli shift of $G$ is sofic by Theorem \ref{theoremabertweiss}, so Bernoulli shifts should come up as a natural example of sofic actions. We can now present the elegant proof of soficity of Bernoulli shifts obtained by Ozawa in full details.% We will not make assumptions on finiteness of the base space of the Bernoulli shifts.

\begin{proposition}[{\cite{MR2566316,ozawasoficnotes,MR2826401}}]
Every Bernoulli shift of a sofic group is sofic.
\end{proposition}
\begin{proof}[Proof (based on \cite{ozawasoficnotes}]
Let $(X,\mu)$ be a (standard) probability space and consider the Bernoulli shift $\beta$ of $G$ on $X^G$. If $G$ is finite, then $G\ltimes X^G$ is periodic and hence sofic (Corollary \ref{periodicequivalencerelationsaresofic}). Assume then that $G$ is infinite so $\beta$ is free.

We assume $G$ is sofic. Let $\theta=(\theta_k)_\mathcal{U}:G\to\prod_\mathcal{U}\mathfrak{S}_{n_k}$ be a sofic embedding. Let $\phi=(\phi_k)_\mathcal{U}:\MAlg(X)\to\prod_\mathcal{U}\MAlg(Y_{m_k})$ be an isometric semigroup embedding (Corollary \ref{periodicequivalencerelationsaresofic}). We may assume that for all $k$, $\phi_k(X)=Y_{m_k}$ and $\phi_k(\varnothing)=\varnothing$.

For each $k$, let $Z_k=Y_{m_k}^{n_k}$ the set of functions $f:Y_{n_k}\to Y_{m_k}$. We identify $Y_{n_k}\times Z_k$ with $Y_{n_k\cdot m_k^{n_k}}$ (since both sets have the same number of elements), and the group $\mathfrak{S}_{Y_{n_k}\times Z_k}$ of permutations on $Y_{n_k}\times Z_k$ with $\mathfrak{S}_{n_k\cdot m_k^{n_k}}$. We will define a covariant sofic embedding $(\psi,\eta)$, given by maps $\eta:G\to\prod_\mathcal{U}\mathfrak{S}_{Y_{n_k}\times Z_k}$ and $\psi:\MAlg(X^G)\to\prod_\mathcal{U}\MAlg(Y_{n_k}\times Z_k)$. First, we define $\psi$ on cylinder sets of $X^G$ (which generate its Borel structure).

Given $A_1,\ldots,A_n\in\MAlg(X)$ and distinct $g_1,\ldots,g_n\in G$, a \emph{cylinder set} is one which can be written as
\[C_{g_1,\ldots,g_n}^{A_1,\ldots,A_n}=\bigcap_{j=1}^n\left\{f\in X^G: f(g_j)\in A_j)\right\}.\]
This representation of a nonempty cylinder set in terms of $g_1,\ldots,g_k$ and $A_1,\ldots,A_k$ is unique up to those $A_j$ which are equal to $X$ (and their order), and such a representation describes the empty set if and only if at least one of the $A_j$ is empty. We write $\mathfrak{C}$ for the collection of cylinder sets in $X^G$. Note that $\mathfrak{C}$ is closed under intersections.

For each $k$, define $\psi_k:\MAlg(X^G)\to\MAlg(Z_k)$ by
\[\psi_k(C_{g_1,\ldots,g_n}^{A_1,\ldots,A_n})=\bigcap_{j=1}^n\left\{(i,f)\in Y_{n_k}\times Z_k:f\left(\theta_k(g_j)^{-1}(i)\right)\in\phi_k(A_j)\right\}\]
By uniqueness of the representation of cylinder sets and since $\phi_k(X)= Y_{m_k}$ and $\phi_k(\varnothing)=\varnothing$, the definition of $\psi$ does not depend on the given choice (and order) of $g_j$ and $A_j$.

Therefore we obtain a map $\psi=(\psi)_\mathcal{U}:\mathfrak{C}\to\prod_\mathcal{U}\MAlg(Y_{n_k}\times Z_k)$, which we later extend to all of $\MAlg(X)$. For this, we first prove that $\psi$ preserves meets (intersections).

%Define $\psi=(\psi_k)_\mathcal{U}$. First let us show that $\psi$ is an isometric embedding of $\MAlg(X)$.

Consider cylinders $C_{g_1,\ldots,g_n}^{A_1,\ldots,A_n}$ and $C_{h_1,\ldots,h_m}^{B_1,\ldots,B_m}$. By modifying the choice and order of $g_i,h_i,A_i,B_i$, we can assume $m=n$ and $g_i=h_i$, so that
\[C_{g_1,\ldots,g_n}^{A_1,\ldots,A_n}\cap C_{g_1,\ldots,g_n}^{B_1,\ldots,B_n}=C_{g_1,\ldots,g_n}^{A_1\cap B_1,\ldots,A_n\cap B_n}.\]
It follows that, for all $k$,
\begin{align*}
\psi_k(C_{g_1,\ldots,g_n}^{A_1,\ldots,A_n})\cap\psi_k(C_{g_1,\ldots,g_n}^{B_1,\ldots,B_n})\hspace{-3cm}&\hspace{3cm}\\
&=\bigcap_{j=1}^n\left\{(i,f)\in Y_{n_k}\times Z_k:f\left(\theta_k(g_j)^{-1}(i)\right)\in\phi_k(A_j)\cap\phi_k(B_j)\right\}\ntag\label{equationbernoullishiftintersection1}
\end{align*}
and
\[\psi_k(C_{g_1,\ldots,g_n}^{A_1,\ldots,A_n}\cap C_{g_1,\ldots,g_n}^{B_1,\ldots,B_n})=\bigcap_{j=1}^n\left\{(i,f)\in Y_{n_k}\times Z_k:f\left(\theta_k(g_j)^{-1}(i)\right)\in\phi_k(A_j\cap B_j)\right\}\ntag\label{equationbernoullishiftintersection2}\]
Let $M_k=\max\left\{d_\#(\phi_k(A_j)\cap\phi_k(B_j),\phi_k(A_j\cap B_j):j=1,\ldots,n\right\}$.

If $(i,f)$ is in the symmetric difference of the sets described in Equations (\ref{equationbernoullishiftintersection1}) and (\ref{equationbernoullishiftintersection2}), then there is $j$ such that
\[f(\theta_k(g_j)^{-1}(i))\in(\phi_k(A_j)\cap\phi_k(B_k))\triangle\phi_k(A_j\cap B_j).\]

For each $j$ and each $i\in Y_{n_k}$, there are exactly
\[m_k^{n_k-1}\cdot\#\big[(\phi_k(A_j)\cap\phi_k(B_k))\triangle\phi_k(A_j\cap B_j)\big]\]
functions from $Y_{n_k}$ to $Y_{m_k}$ which map $\theta_k(g_j)(i)$ to an element of
\[(\phi_k(A_j)\cap\phi_k(B_k))\triangle\phi_k(A_j\cap B_j).\]
Letting $j\in\left\{1,\ldots,n\right\}$ and $i\in Y_{n_k}$ go through all possible values, we conclude that the (normalized Hamming) distance of the sets in Equations (\ref{equationbernoullishiftintersection1}) and (\ref{equationbernoullishiftintersection2}) is at most
\begin{align*}
\frac{1}{n_km_k^{n_k}}\cdot n_k\cdot \sum_{j=1}^nm_k^{n_k-1}\#\big[(\phi_k(A_j)\cap\phi_k(B_k))\triangle\phi_k(A_j\cap B_j)\big]\hspace{-3cm}&\hspace{3cm}=\\
&\sum_{j=1}^nd_{\#}(\phi_k(A_j)\cap\phi_k(B_j),\phi_k(A_j,B_j))
\end{align*}
which converges to $0$ along $\mathcal{U}$, because $\phi=(\phi_k)_\mathcal{U}$ is a sofic embedding. Therefore, $\psi$ satisfies $\psi(C_1\cap C_2)=\psi(C_1)\land\psi(C_2)$ for all cylinders $C_1$, $C_2$, and moreover $\psi(\varnothing)=0$. This allows us to extend $\psi$ to a semigroup embedding to the (Boolean) algebra of sets generated by cylinders, by setting
\[\psi(A)=\bigvee_k\psi(C_k,)\]
where $\left\{C_k\right\}_k$ is any finite collection of pairwise disjoint cylinders whose union is $A$. This way, $\psi$ still preserves joins and meets.

Now we show $\psi$ is isometric, or equivalently that it preserves traces (measures)  by Theorem \ref{theoremtracepreservingisisometric}. But again, since $\psi$ preserves joins, it is enough to prove that $\psi$ preserves the trace of cylinders, so consider a cylinder of the form
\[C=C_{g_1,\ldots,g_n}^{A_1,\ldots,A_n}.\]
For all $k\in\mathbb{N}$, let
\[I_k=\left\{i\in Y_{n_k}:\theta_{g_1}(i),\theta_{g_2}(i),\ldots,\theta_{g_n}(i)\text{ are pairwise distinct}\right\}.\]
Since $\theta$ is a sofic embedding, $\mu_\#(I_k)$ converges to $1$ along $\mathcal{U}$.

If $i\in I_k$, then there are exactly
\[m_k^{n_k-n}\cdot\prod_{j=1}^n\#\left(\phi_k(A_j)\right)\]
functions $f:Y_{n_k}\to Y_{m_k}$ for which $(i,f)\in\psi_k(C)$. Therefore,

\[\begin{array}{rcccc}
\mu_\#(\psi_k(C))&=&\mu_\#(\psi_k(C)\cap(I_k\times Z_k))&+&\mu_\#(\psi_k(C)\setminus(I_k\times Z_k))\vspace{1ex}\\
&\leq&\displaystyle{\frac{\#(I_k)}{n_km_k^{n_k}}m_k^{n_k-n}\cdot\prod_{j=1}^n\left(\#\phi_k(A_j)\right)}&+&\displaystyle{\mu_\#(\psi_k(C)\setminus(I_k\times Z_k))}\vspace{1ex}\\
&=&\displaystyle{\mu_\#(I_k)\prod_{j=1}^n\mu_\#(A_j)}&+&\displaystyle{\mu_\#(\psi_k(C)\setminus(I_k\times Z_k)).}
\end{array}\]
The first term above converges to $\prod_{j=1}^n\mu_{\#}(A_j)$, which is exactly the measure of $C$, whereas the second term is bounded by
\[\mu_\#(\psi_k(C)\setminus(I_k\times Z_k))\leq\mu_\#((Y_{n_k}\times Z_k)\setminus(I_k\times Z_k))=\mu_\#(Y_{n_k}\setminus I_k)=1-\mu_\#(I_k)\]
which converges to $0$ along $\mathcal{U}$. Therefore $\psi$ is isometric, and since the algebra generated by cylinders is dense in $\MAlg(X^G)$, $\psi$ can be extended to an isometric semigroup embedding $\psi:\MAlg(X^G)\to\prod_{\mathcal{U}}\MAlg(Y_{n_k}\times Z_k)$.

Now we let, for each $k$, $\eta_k:G\to\mathfrak{S}_{Y_{n_k}\times Z_k}$ be defined by
\[\eta_k(g)(i,f)=(\theta_g(i),f),\qquad\text{for all }g\in G\quad\text{and}\quad(i,f)\in Y_{n_k}\times Z_k.\]
Clearly,
%\tr_\#(\eta_k(g))=\tr_\#(\theta_k(g))\quad\text{and}\quad 
\[d_\#(\eta_k(g)\eta_k(h),\eta_k(gh))=d_{\#}(\theta_k(g)\theta_k(h),\theta_k(gh))\]
for all $g,h\in G$. Thus we define a morphism $\eta=(\eta_k)_\mathcal{U}:G\to\prod_\mathcal{U}\mathfrak{S}_{Y_{n_k}\times Z_k}$.

We denote by $\beta$ the Bernoulli shift of $G$ on $X$. It remains only to show that $(\eta,\psi)$ is covariant, and again it is enough to prove this for cylinder sets. Let $g\in G$ and a cylinder $C_{g_1,\ldots,g_n}^{A_1,\ldots,A_n}$. Then
\begin{itemize}
\item $\displaystyle{\beta_gC_{g_1,\ldots,g_n}^{A_1,\ldots,A_n}=C_{gg_1,\ldots,gg_n}^{A_1,\ldots,A_n}}$;
\item $\displaystyle{\eta_k(g)(\psi_k(C_{g_1,\ldots,g_n}^{A_1,\ldots,A_n}))=\bigcap_{j=1}^n\left\{(i,f):f(\theta_k(g_j)^{-1}\theta_k(g)^{-1}(i))\in\phi_k(A_j)\right\}}$; and
\item $\displaystyle{\psi_k(C_{gg_1,\ldots,gg_n}^{A_1,\ldots,A_n})=\bigcap_{j=1}^n\left\{(i,f):f(\theta_k(gg_j)^{-1}(i))\in\phi_k(A_j)\right\}}$.
\end{itemize}
If $(i,f)$ is in the symmetric difference $\eta_k(g)(\psi_k(C_{g_1,\ldots,g_n}^{A_1,\ldots,A_n}))\triangle\psi_k(C_{gg_1,\ldots,gg_n}^{A_1,\ldots,A_n})$, then there is $j$ for which
\[\theta_k(gg_j)^{-1}(i)\neq\theta_k(g_j)^{-1}\theta_k(g)^{-1}(i)\]
and there are $m_k^{n_k}$ choices for $f\in Z_k$. Thus the distance between these two sets is at most
\[\frac{1}{n_k m_k^{n_k}}\sum_{j=1}^nn_kd_\#(\theta_k(gg_j)^{-1},\theta_k(g_j)^{-1}\theta_k(g)^{-1})\cdot m_k^{n_k}=\sum_{j=1}^nd_\#(\theta_k(gg_j)^{-1},\theta_k(g_j)^{-1}\theta_k(g)^{-1})\]
which converges to $0$ since $\theta$ is a sofic embedding. This proves that $\eta(g)(\psi(C))=\psi(g\cdot C)$ for all $g\in G$ and all cylinders $C$, and hence for all $C\in\MAlg(X^G)$. Therefore $\beta$ is sofic by Theorem \ref{theoremrewriesoficactionintermsofcovariant}.\qedhere
\end{proof}

%\begin{corollary}
%A group is sofic if and only if it admits a (essentially free) sofic action.
%\end{corollary}

One important outstanding question in the area is whether all actions of sofic groups are sofic. However, we may study the related question of how large the class of groups with sofic actions is, by proving that it is closed under some group-theoretic constructions. First, let us quickly review, without proofs, the known closure properties of the class of sofic groups.

%\begin{lemma}
%A group $G$ is sofic if and only if every finitely generated subgroup of $G$ is sofic.
%\end{lemma}

\begin{theorem}\label{theoremclosednessofsofic}
The class of sofic groups is closed under:
\begin{enumerate}
\item[(a)] Products \cite{MR2220572}: If $\left\{G_i\right\}_{i\in I}$ is any family of sofic groups, then $\prod_{i\in I}G_i$ is sofic;
\item[(b)] Subgroups \cite{MR2220572}: If $G$ is sofic and $H\subseteq G$ is a subgroup, then $H$ is sofic;
\item[(c)] Direct limits \cite{MR2220572}: If $\left(\phi_{j,i}: G_i\to G_j\right)_{i,j}$ is a direct net of sofic groups, then the direct limit $\lim_{i\to\infty}G_i$ is sofic;
\item[(d)] Amenable-by-sofic extensions \cite{MR2220572}: If $N$ is a normal subgroup of $G$, $G/N$ is amenable and $N$ is sofic, then $G$ is sofic;
\item[(e)] Unrestricted wreath products by amenable groups \cite{arxiv1802.04688v1}: If $G$ is sofic and $A$ is amenable then $G\wr\!\wr A$ is amenable
\item[(f)] Free products amalgamated over an amenable subgroup \cite{MR2826401,MR2823074}: If $G$ and $H$ are sofic and $A$ is a common amenable subgroup of $G$ and $H$, then $G\ast_A H$ is sofic.
\item[(g)] Wreath products \cite{MR3795485}: If $G$ and $H$ are sofic then $G\wr H$ is sofic;
\item[(h)] Graph limits (see \cite{MR2794910} for details).
\end{enumerate}
\end{theorem}

\begin{definition}[{\cite{MR2826401}}]
We denote by $\mathscr{S}$ the class of groups whose probability measure-preserving actions are all sofic.
\end{definition}

We will be interested in proving that the class $\mathscr{S}$ is closed under some of the constructions listed in \ref{theoremclosednessofsofic}. %Although at the moment we cannot prove whether the class of \emph{elementary sofic} group, defined in \cite{MR2794910} by Cournulier as the smallest class of groups closed under subgroups, direct limits, amenable extensions and graph limits.

\begin{theorem}[{\cite{MR2826401}}]
A group $G$ belongs to the class $\mathscr{S}$ if and only if every free probability-measure preserving action of $G$ is sofic.
\end{theorem}
\begin{proof}
Suppose that every free probability-measure preserving action of $G$ is sofic. Let $\alpha$ be any probability measure-preserving action of $G$ on $(X,\mu)$. Let $\beta$ be any free probability-measure preserving action of $G$ on a space $(Y,\mu)$ (e.g. a Bernoulli shift on a non-atomic space), and define a new free probability measure-preserving action $\Gamma$ of $G$ on $X\times Y$ by
\[\Gamma_g(x,y)=(\alpha_g(x),\beta_g(y))\]
by assumption, $\Gamma$ is sofic, so there exists a sofic embedding $\Theta:\MEAS(G\ltimes_\Gamma (X\times Y))\to\prod_\mathcal{U}\MEAS(Y_{n_k}^2)$.

We can embed $\MEAS(G\ltimes_\alpha X)$ isometrically into $\MEAS(G\ltimes_\Gamma (X\times Y))$, via $A\mapsto A\times Y$, so composing this with $\Phi$ we obtain a sofic embedding of $G\ltimes_\alpha X$.\qedhere
\end{proof}

%\begin{lemma}
%A group $G$ belongs to $\mathscr{S}$ if and only if every finitely generated subgroup of $G$ belongs to $\mathscr{S}$.
%\end{lemma}

\begin{theorem}
The class $\mathscr{S}$ is closed under
\begin{enumerate}
\item[(a)] Subgroups: If $G\in\mathscr{S}$ and $H\subseteq G$ is a subgroup, then $H\in\mathscr{S}$;
%\item[(b)] Products: If $\left\{G_i\right\}_{i\in I}$ is any family of sofic groups, then $\prod_{i\in I}G_i$ is sofic;
\item[(b)] Direct limits: If $\left(\phi_{j,i}: G_i\to G_j\right)$ is a direct net of groups in $\mathscr{S}$, then $\lim_{i\to\infty}G_i\in\mathscr{S}$;
\item[(c)] Co-amenable extensions: If $N$ is a co-amenable subgroup of $G$ (that is, the left action of $G$ on $G/N$ is amenable) and $N\in\mathscr{S}$, then $G\in\mathscr{S}$;
%\item[(e)] Wreath products \cite{hayes_sale_metric}: If $G$ and $H$ are sofic then $G\wr H$ is sofic;
%\item[(f)] Unrestricted wreath products by amenable groups \cite{arxiv1802.04688v1}: If $G$ is sofic and $A$ is amenable then $G\wr\wr A$ is amenable
\item[(d)] Free products amalgamated over an amenable subgroup \cite[Theorem 3.17]{MR2826401}: If $G,H\in\mathscr{S}$ and $A$ is a common amenable subgroup of $G$ and $H$, then $G\ast_A H$ is sofic.
%\item[(h)] Graph limits; If $N,N_n$ are normal subgroups of $\mathbb{F}$ and such that $\mathbb{F}/N_n\in\mathscr{S}$ and $N_n$ converges to $N$, then $\mathbb{F}/N\in\mathscr{S}$.
\end{enumerate}
\end{theorem}
\begin{proof}
\begin{enumerate}
\item[(a)] Suppose $G\in\mathscr{S}$ and $H$ is a subgroup of $G$. We let $\mathbb{N}_0=\mathbb{N}\cup\left\{0\right\}$ the set of non-negative integers. To simplify the notation, let us assume that the index $[G:H]$ is infinite, although the same arguments work in the case of finite index.

Let $\left\{a_i:i\in\mathbb{N}_0\right\}$ be a complete set of left representatives of $H$ in $G$ with $a_0=1$.
For each pair $(g,i)\in G\times \mathbb{N}_0$, let $J(g,i)\in\mathbb{N}_0$ be the unique number such that $ga_i\in a_{J(g,i)} H$. Then for all $g$, $a_i^{-1}g^{-1}a_{J(g,i)}\in H$

Note that, if we identify $\mathbb{N}_0$ with the quotient space $G/H$ with via $i\mapsto a_iH$, the map $J:G\times\mathbb{N}_0\to\mathbb{N}_0$ corresponds to the action map of the usual left action of $G$ on $G/H$.

Let us show that any free probability measure-preserving action $\alpha$ of $H$ on a space $(X,\mu)$ is sofic. First, we construct a probability measure-preserving action $\theta$ of $G$ on $X^{\mathbb{N}_0}$ by setting
\[\left[\theta_g({\mathbf{x}})\right]_i=\alpha(a_i^{-1}ga_{J(g^{-1},i)})(x_{J(g^{-1},i)}),\qquad\text{for all }{\mathbf{x}}=(x_i)_i\in X^{\mathbb{N}_0}\]
Then $\theta$ is a measure-preserving action of $G$ on $X^{\mathbb{N}_0}$, so we can take a sofic embedding $\Theta:\MEAS(G\ltimes_\theta X^{\mathbb{N}_0})\to\prod_\mathcal{U}\MEAS(Y_n^2)$. The map
\[\iota:\MEAS(H\ltimes_\alpha X)\to\MEAS(G\ltimes_\theta X^{\mathbb{N}_0})\]
\[\iota(A)=\left\{(g,(x_i)_i)\in G\times X^{\mathbb{N}_0}:(g,x_0)\in A\right\}\]
is an isometric semigroup morphism, because $\alpha_0=1$, so $\Theta\circ\iota$ is a sofic embedding for $H\ltimes_\alpha X$.
%\item[(b)] We first prove that $\mathscr{S}$ is closed under finite products. Suppose $G,H\in\mathscr{S}$, and $\theta$ is a probability measure-preserving action of $G\times H$ on a space $(X,\mu)$. Restricting $\theta$ to $G$ and $H$, we obtain probability-measure preserving actions, so take covariant sofic embeddings $(\sigma_G,\phi_G)$ and $(\sigma_H,\phi_H)$. We may assume $\phi_G=\phi_H=\phi$ by \ref{}. Let $\sigma=\sigma_G\ast\sigma_H:G\ast H\to\prod_\mathcal{U}\mathfrak{S}_{n_k}$ be the free product of $\sigma_G$ and $\sigma_H$. Then $(\sigma,\phi)$ is covariant for the action of $G\ast H$ on $X$. By Lemma \ref{},
\item[(b)] Suppose $G=\displaystyle{\lim_{\longrightarrow}}G_i$, where $G_i\in\mathscr{S}$. For each $i$, let $p_i:G_i\to G$ be the canonical morphism. Let $\alpha$ be a free probability measure-preserving action of $G$ on $(X,\mu)$.

Given a finite $F\subseteq G$, choose $i$ such that $F\cup F^2\subseteq p_i(G_i)$, and for each $x\in F\cup F^2$ choose $x^i\in G_i$ such that $p_i(x^i)=x$.

Given $x,y\in F$, we do not necessarily have $(x^i)(y^i)=(xy)^i$, but since
\[p_i((xy)^i)=xy=p_i(x^iy^i)\]
then there is $j\geq i$ such that $p_{j,i}(x^iy^i)=p_{j,i}((xy)^i)$, so substituting $i$ by $j$ and $x^i$ by $p_{j,i}(x^i)$ for all $x\in F\cup F^2$ if necessary, we may assume $x^iy^i=(xy)^i$ for all $x\in F$.

Let $\alpha_i$ be the composition of $\alpha$ with $p_i$, which is a probability measure-preserving but not necessarily free. In any case, $\alpha_i$ is sofic, so we can find a covariant pair of morphisms
\[\eta:G_i\to\prod_\mathcal{U}\mathfrak{S}_{n_k}\quad\text{and}\quad\phi:\MAlg(X)\to\prod_\mathcal{U}\MAlg(Y_{n_k}),\]
and defining $\sigma(f)=\eta(f_i)$ we obtain, for all $f\in F$ and $A\in\MAlg(A)$,
\[\phi(f\cdot A)=\phi(f_i\cdot A)=\eta(f_i)\cdot\phi(A)=\sigma(f)\cdot\phi(A)\]
and
\[\sigma(fg)=\eta((fg)^i)=\eta(f^i)\eta(g^i)=\sigma(f)\sigma(g)\]
so using Theorem \ref{theoremequivalenceembeddingintoultraproductformetricstructures} we can obtain almost covariant morphisms for $F$ (and any finite subset of $\MAlg(X)$), in the sense of Corollary \ref{corollaryalmostcovariant}. This proves that $\alpha$ is sofic.
\item[(c)] Let $\alpha$ be a free probability measure-preserving action of $G$ on $(X,\mu)$, $F\subseteq G$ finite, $\mathcal{A}\subseteq\MAlg(X)$ finite and $\epsilon>0$. Using co-amenability of $N\subseteq G$, choose a finite subset $P\subseteq G/N$ such that
\[\#((F\cdot P)\triangle P)\leq\epsilon\cdot\#(P)\]
where $F\cdot P=\left\{(fg)N:f\in F, gN\in P\right\}$. Let $a\mapsto \overline{a}$ be a section of the quotient map $G\to G/N$ -- i.e., $\overline{a}N=a$ for all $a\in G/N$.

The restriction $\alpha|_N$ is free and probability measure-preserving, and for all $f\in F$ and $a\in P$,
\[(\overline{ga})^{-1}g\overline{a}\in N.\]
Let $F'=\left\{(\overline{ga})^{-1}g\overline{a}:g\in F^2\cup F, a\in P\right\}$, which is a finite set. Similarly, let $\mathcal{A}'=\left\{\alpha(\overline{a}^{-1})(A):a\in P, A\in\mathcal{A}\right\}$, which is also a finite subset of $\MAlg(X)$.

Since $N\in\mathscr{S}$ we can find $n\in\mathbb{N}$ and maps
\[\sigma:N\to\mathfrak{S}_n\quad\text{and}\quad\phi:\MAlg(X)\to\MAlg(Y_n)\]
which form a $(F',\mathcal{A}',\epsilon)$-almost covariant morphism. Define\ftnt{Again, we denote by $\mathfrak{S}_Z$ the permutation group of a set $Z$.} $\pi:G\to\mathfrak{S}_{Y_n\times P}$ by setting
\[\pi(g)(i,a)=(\sigma((\overline{ga})^{-1}g\overline{a})(i),ga),\]
which is defined on a set of measure $\geq 1-\epsilon$. Then for all $g,h\in F$ we have
\begin{align*}
\pi(g)\pi(h)(i,a)&\overset{\epsilon}{=}\pi(g)(\sigma((\overline{ha})^{-1}h\overline{a})(i),ha)\\
&\overset{\epsilon}{=}(\sigma((\overline{gha})^{-1}g\overline{ha})\sigma((\overline{ha})^{-1}h\overline{a})(i),gha)\\
&\overset{\epsilon}{=}(\sigma((\overline{gha})^{-1}g\overline{ha}(\overline{ha})^{-1}h\overline{a})(i),gha)\\
&=(\sigma((\overline{gha})^{-1}gh\overline{a})(i),gha)\overset{\epsilon}{=}\pi(gh)(i,a)
\end{align*}
where the ``$\overset{\epsilon}{=}$'' means that this equality is valid only up to a set of measure $\epsilon$. So all these terms are equal on a set of measure $\geq1-5\epsilon$, which means that $d_\#(\pi(g)\pi(h),\pi(gh))\leq5\epsilon$. This corresponds to property (i) of an almost covariant morphism (Corollary \ref{corollaryalmostcovariant}).

We also define $\psi:\MAlg(X)\to\MAlg(Y_n\times P)$ by setting
\[\psi(A)=\bigcup_{a\in P}\phi(\alpha(\overline{a}^{-1})(A))\times\left\{a\right\}\]
and easy computations, similar to the one above, prove that properties (ii) and (iii) of Corollary \ref{corollaryalmostcovariant} are satisfied. We prove almost covariance of $\pi$ and $\psi$. Using the notation ``$\overset{\epsilon}{=}$'' as above we have, for $g\in F$ and $A\in\mathcal{A}$,
\begin{align*}
\pi(g)(\psi(A))&\overset{\epsilon}{=}\bigcup_{a\in P}\sigma(\overline{ga}^{-1}g\overline{a})\left[\phi(\alpha(\overline{a}^{-1}))(A)\right]\times\left\{ga\right\}\\
&\overset{\epsilon}{=}\bigcup_{a\in P}\phi(\alpha(\overline{ga}^{-1}g)(A))\times\left\{ga\right\}\\
&\overset{\epsilon}{=}\bigcup_{b\in P}\phi(\alpha(\overline{b}^{-1}g)(A))\times\left\{b\right\}\\
&\overset{}{=}\bigcup_{b\in P}\phi(\alpha(\overline{b}^{-1})(\alpha(g)(A)))\times\left\{b\right\}\\
&=\psi(\alpha(g)(A))
\end{align*}
and therefore $(\pi,\psi)$ is a $(F,\mathcal{A},5\epsilon)$-almost covariant morphism.
\item[(d)] We refer to \cite[Theorem 3.17]{MR2826401}.\qedhere
\end{enumerate}
\end{proof}
%\end{document}